% The paper (PLAN Phase 6). Every number and every results sentence comes from generated/, written by
% python -m ecgaug.paper (and the final evaluation) from logged runs; never type a number by hand. A red
% [bracket] is a result that does not exist yet. Citations: refs.bib holds only verified entries.
\documentclass[11pt]{article}
\usepackage[margin=2.5cm]{geometry}
\usepackage{booktabs}
\usepackage{graphicx}
\usepackage{adjustbox}
\usepackage{amsmath}
\usepackage{xcolor}
\usepackage[colorinlistoftodos]{todonotes}
\usepackage{hyperref}
\hypersetup{hidelinks}
\newcommand{\missing}[1]{\textcolor{red}{[#1]}}  % a result that does not exist yet (written by python -m ecgaug.paper)
% Generated by python -m ecgaug.paper; do not edit. Every value comes from a result file.
\newcommand{\CountTestF}{388}
\newcommand{\CountTestN}{44,259}
\newcommand{\CountTestS}{1,837}
\newcommand{\CountTestV}{3,221}
\newcommand{\CountTrainingF}{388}
\newcommand{\CountTrainingN}{35,496}
\newcommand{\CountTrainingS}{644}
\newcommand{\CountTrainingV}{3,030}
\newcommand{\CountValidationF}{27}
\newcommand{\CountValidationN}{10,370}
\newcommand{\CountValidationS}{300}
\newcommand{\CountValidationV}{758}
\newcommand{\PhaseOneInputRepresentationAMean}{0.164}
\newcommand{\PhaseOneInputRepresentationBMean}{0.001}
\newcommand{\PhaseOneInputRepresentationCMean}{0.296}
\newcommand{\PhaseOneInputRepresentationDMean}{0.035}
\newcommand{\PhaseOneWinner}{C}


\title{Does Generative Augmentation Help Rare-Arrhythmia Detection Under Inter-Patient Evaluation?\\
A Controlled Comparison of Conditional VAEs and WGAN-GP}
\author{\missing{authors}}
\date{}

\begin{document}
\maketitle

\begin{abstract}
Class-conditional generative models are often proposed to rebalance electrocardiogram (ECG) beat classifiers
toward rare arrhythmias, but many reported gains come from evaluations in which beats of the same patient
appear in training and test. We ask whether generative augmentation helps when that leak is closed. On the
MIT-BIH Arrhythmia Database with the inter-patient DS1/DS2 protocol, we compare a conditional variational
autoencoder (cVAE) and a Wasserstein GAN with gradient penalty (WGAN-GP) against no augmentation,
class-weighted loss, random oversampling and SMOTE, with two classifier architectures and three seeds. Every
model, strength and checkpoint is chosen on validation by rules fixed before the experiments, the test set is
read once, and uncertainty is estimated by resampling test patients. \missing{results sentence: written after /final-eval}

We also measure how much intra-patient evaluation inflates results, test whether synthetic-quality metrics
predict the downstream benefit, examine whether augmentation changes what the classifiers rely on, and
generate grounded clinician-readable notes whose every claim is checked against the model output.
\end{abstract}

\section{Introduction}
Automatic beat classification on the MIT-BIH Arrhythmia Database~\cite{moody2001,goldberger2000} is one of the
oldest benchmarks in biomedical signal processing. The clinically important classes are rare:
supraventricular ectopic beats (S) and fusion beats (F) together make up a few percent of the beats, and
they are concentrated in a handful of patients. Classifiers trained on the imbalanced data favour normal
beats, and a large literature proposes to rebalance them with synthetic beats from generative adversarial
networks or variational autoencoders~\cite{delaney2019ecggan,shaker2020gan,adib2022cgan}.

Two problems make these results hard to read. First, many studies split beats at random, so beats of the
same patient appear in training and test; a classifier can then recognise the patient rather than the
arrhythmia. The inter-patient protocol of de Chazal et al.~\cite{dechazal2004}, which trains on one set of
records (DS1) and tests on another (DS2), removes this leak, and results under it are much lower; a recent
systematic review finds that only a minority of studies follow it~\cite{silva2025review}. Second, the
choices that decide the outcome (architecture, augmentation strength, training epoch, the generator
checkpoint) are rarely fixed before the test set is consulted.

This paper answers a narrow question carefully: under a strict inter-patient protocol, does
class-conditional generative augmentation improve detection of rare arrhythmias compared with simple
rebalancing, and what explains when it helps? Our contributions are:
\begin{itemize}
\item \textbf{A controlled, leakage-free comparison} of cVAE and WGAN-GP augmentation against no augmentation,
class-weighted loss, random oversampling and SMOTE, with two architectures, several seeds, record-level
uncertainty and selection rules registered before the experiments they govern.
\item \textbf{A leakage audit}: the same pipeline under a random beat-level split, to measure how much
intra-patient evaluation inflates the baseline and the apparent benefit of augmentation.
\item \textbf{A test of synthetic-quality metrics}: whether DTW, spectral, FID-analog, TSTR, diversity and
two-sample-test metrics predict the downstream gain of a generator checkpoint.
\item \textbf{An explainability lens}: whether augmentation changes the faithfulness and physiological alignment
of saliency maps, including reliance on artifacts of synthetic beats.
\item \textbf{A grounded reporting layer}: clinician-readable notes generated from the model output, with a fact
checker that rejects any unsupported claim.
\end{itemize}
Negative results are part of the design: every pre-registered comparison is reported, whatever its outcome.

\section{Related work}
\paragraph{Evaluation protocols.} The AAMI EC57 standard~\cite{aami2012ec57} groups MIT-BIH beat labels into
the classes N, S, V, F and Q and asks for record-by-record reporting. De Chazal et al.~\cite{dechazal2004}
introduced the DS1/DS2 inter-patient split used here; Llamedo and Mart\'inez~\cite{llamedo2011} selected
features by how well they generalise across databases. Silva et
al.~\cite{silva2025review} review recent work against these standards.

\paragraph{Generative augmentation of ECG.} Recurrent conditional GANs were proposed for medical time
series and evaluated by training on synthetic and testing on real data~\cite{esteban2017rcgan}. GANs have
been used to synthesise realistic ECG~\cite{delaney2019ecggan} and to augment MIT-BIH classifiers, where they were
compared with other augmentation techniques~\cite{shaker2020gan}; an auxiliary-classifier WGAN-GP with DTW-based screening is
closest to our WGAN-GP variants~\cite{adib2022cgan}. Diffusion models with structured state-space backbones
now generate conditional 12-lead ECG~\cite{alcaraz2023sssd}; we leave this family to future work. Our
generators build on the VAE~\cite{kingma2013vae}, its conditional form~\cite{sohn2015cvae}, the Wasserstein
GAN~\cite{arjovsky2017wgan} with gradient penalty~\cite{gulrajani2017wgangp}, and the conditioning schemes of
auxiliary classifiers~\cite{odena2016acgan}, projection critics~\cite{miyato2018projection} and
FiLM~\cite{perez2018film}. Oversampling baselines are random oversampling and SMOTE~\cite{chawla2002smote}.

\paragraph{Synthetic-quality metrics.} We score synthetic beats with dynamic time warping~\cite{sakoe1978dtw},
spectral distances, a Fr\'echet distance in the feature space of a fixed network after
FID~\cite{heusel2017fid}, train-on-synthetic-test-on-real~\cite{esteban2017rcgan} and classifier two-sample
tests~\cite{lopezpaz2017c2st}.

\paragraph{Explanations and reports.} We compare Grad-CAM~\cite{selvaraju2016gradcam}, integrated
gradients~\cite{sundararajan2017ig}, SmoothGrad~\cite{smilkov2017smoothgrad} and occlusion. Large ECG-language
models can write diagnostic reports~\cite{zhao2024ecgchat}; our note generator instead restates the classifier
output and is checked claim by claim.

\section{Data and protocol}
We use the MIT-BIH Arrhythmia Database~\cite{moody2001} from PhysioNet~\cite{goldberger2000}: 48 half-hour
two-lead recordings at 360\,Hz. Lead MLII is selected by name. The four paced records are excluded. Reference
beat annotations are mapped to the AAMI classes N, S, V and F; the few Q beats are excluded from training and
metrics, and RR intervals are measured between consecutive QRS annotations before any class is dropped. A beat
is usable when its window fits in the record and both neighbouring RR intervals exist; the same rule holds for
every representation.

DS1 (22 records) provides training and validation; DS2 (22 records) is the test set. The validation set is a
patient-wise subset of DS1 (records 115, 118, 201, 205 and 223), chosen from DS1 counts only. Because it holds
fewer F beats than the pre-registered minimum support, selection uses S F1 with V F1 as tie-break; the test
metric is unchanged. DS1 record 201 and DS2 record 202 come from the same subject; we keep the standard split
for comparability and repeat every test analysis without record 202. The test arrays are readable only by the
final evaluation, and every access to DS2 is logged. Table~\ref{tab:data} gives the counts.
\begin{table}[htbp]\centering\small% Generated by python -m ecgaug.paper from: results/phase1_data/record_counts.csv
\begin{tabular}{lrrrrr}
\toprule
Split & Records & N & S & V & F \\
\midrule
training & 17 & 35,496 & 644 & 3,030 & 388 \\
validation & 5 & 10,370 & 300 & 758 & 27 \\
test (DS2) & 22 & 44,259 & 1,837 & 3,221 & 388 \\
\bottomrule
\end{tabular}
\caption{Records and beats per split.}\label{tab:data}\end{table}

\subsection{Input representation}
Four representations were compared with a small probe CNN on validation (Table~\ref{tab:repr}): a 0.5\,s window
centred on the R peak (180 samples), a 1\,s band-passed window, the 0.5\,s window plus three RR features
(previous and next RR interval, and the previous RR over the mean of the ten before), and the 0.5\,s window
resampled to 125\,Hz. Beats are z-scored one by one. The rule chose variant \PhaseOneWinner, and the dataset was
frozen with SHA-256 hashes before any later experiment.
\begin{table}[htbp]\centering\small\adjustbox{max width=\linewidth}{% Generated by python -m ecgaug.paper from: results/phase1_input_representation/A.json, results/phase1_input_representation/B.json, results/phase1_input_representation/C.json, results/phase1_input_representation/D.json
\begin{tabular}{llrrrr}
\toprule
Variant & Name & S F1 (val.) & V F1 & F F1 & N recall \\
\midrule
A & short window & 0.164 $\pm$ 0.078 & 0.288 & 0.006 & 0.823 \\
B & long window bandpass & 0.001 $\pm$ 0.002 & 0.449 & 0.000 & 0.991 \\
C & short window rr features & 0.296 $\pm$ 0.038 & 0.672 & 0.029 & 0.973 \\
D & short window 125hz & 0.035 $\pm$ 0.007 & 0.326 & 0.039 & 0.688 \\
\bottomrule
\end{tabular}
}\caption{Input representations (validation, probe CNN).}\label{tab:repr}\end{table}

\section{Methods}
\subsection{Classifiers}
Four baselines were trained (Table~\ref{tab:baselines}): a shallow CNN (three convolution blocks, global average
pooling), a 1D ResNet, a two-layer BiLSTM and a CNN--BiLSTM hybrid; RR features enter through a small side
branch. The rule chose one convolutional (\emph{cnn}) and one recurrent (\emph{rnn}) architecture. All
classifiers use unweighted cross-entropy (except the class-weighted condition), AdamW with weight decay 0.01,
batches of 256, learning rate $10^{-3}$ or $3\cdot10^{-4}$ (the better on validation), at most 40 epochs, and
early stopping with patience 5 on the validation selection metric, never before epoch 10. Decisions are argmax;
no threshold is tuned.
\begin{table}[htbp]\centering\small\adjustbox{max width=\linewidth}{% Generated by python -m ecgaug.paper from: nothing yet
\begin{tabular}{l}
\toprule
Architecture \\
\midrule
\missing{phase2\_baseline\_classifiers results} \\
\bottomrule
\end{tabular}
}\caption{Baseline classifiers (validation).}\label{tab:baselines}\end{table}

\subsection{Synthetic-quality metrics}
Four metric families score a set of synthetic beats against real training beats: DTW distance to the five
nearest real beats (Sakoe--Chiba band of 10\%), the distance between Welch spectra, a Fr\'echet distance between
features of a fixed probe network, and TSTR (the validation score of a probe CNN trained on synthetic beats).
Before use, each family had to pass a sanity harness: on halves of the real training beats, its score had to
worsen monotonically under four degradations (noise, label shuffling, time shift, smoothing) at three
severities (Spearman $\rho \ge 0.8$) and vary little across resamples (median CV below 10\%)
(Table~\ref{tab:quality}). Three diagnostics guard against degenerate generators: the diversity ratio (spread of
synthetic beats relative to real ones), the memorization fraction (synthetic beats closer to a training beat
than the 5th percentile of leave-one-out distances among real beats) and the accuracy of a classifier two-sample
test~\cite{lopezpaz2017c2st}.
\begin{table}[htbp]\centering\small\adjustbox{max width=\linewidth}{% Generated by python -m ecgaug.paper from: nothing yet
\begin{tabular}{l}
\toprule
Metric \\
\midrule
\missing{phase2\_quality\_metrics results} \\
\bottomrule
\end{tabular}
}\caption{Sanity tests of the quality metrics.}\label{tab:quality}\end{table}

\subsection{Generators}
Both families learn beats of all four classes with class-balanced batches and generate S, V and F beats; the RR
features of a synthetic beat are copied from a random real training beat of the same class. Four cVAE variants
(convolutional with concatenated class input; the same with a smaller KL weight; with an added spectral loss;
a residual network with FiLM conditioning~\cite{perez2018film}) and four WGAN-GP variants (concatenated class
embedding; projection critic~\cite{miyato2018projection}; auxiliary classifier~\cite{odena2016acgan}; one
unconditional GAN per class) were piloted with equal step budgets on two seeds. At each quarter of the budget a
pilot classifier (the chosen cnn) was trained on real plus synthetic beats; hard gates (diversity ratio at least
0.7, memorization fraction at most 10\%) were applied to every class and seed, and the variant with the best
validation score at the final checkpoint was chosen per family (Table~\ref{tab:generators}). The chosen variants
were trained at the full budget on every grid seed, and a synthetic bank was exported from the final checkpoint
and frozen with its hash. Generators see training records only.
\begin{table}[htbp]\centering\small\adjustbox{max width=\linewidth}{% Generated by python -m ecgaug.paper from: nothing yet
\begin{tabular}{llllrrrl}
\toprule
Family & Variant & Setting & Pilot S F1 & Diversity (min) & Memorization (max) & C2ST (max) & Hard gates \\
\midrule
\missing{phase 3 pilot results} &  &  &  &  &  &  &  \\
\bottomrule
\end{tabular}
}\caption{Generator pilots: pilot score, diagnostics and hard gates.}\label{tab:generators}\end{table}

\subsection{Augmentation policy and comparison grid}
Four policies decide how many synthetic beats each minority class receives: balancing to the training N count,
a fraction of the N count, a multiple of the class's real count, and a fraction of the N count drawn from a
DTW-screened pool. The policy was chosen on validation with both banks, subject to a limit of 0.01 on the loss of
N recall (Table~\ref{tab:policy}). The full grid then trains six conditions (none, class-weighted loss with
weights $(n_{\max}/n_c)^{\gamma}$, random oversampling, SMOTE, cVAE, WGAN-GP) at three strengths each (one for
none), for both architectures and three seeds; classifier seed $s$ uses bank seed $s$, and validation and test
stay entirely real. The strength of each condition and the best generative condition per architecture were
fixed on validation before the test set was read.
\begin{table}[htbp]\centering\small\adjustbox{max width=\linewidth}{% Generated by python -m ecgaug.paper from: nothing yet
\begin{tabular}{l}
\toprule
Policy \\
\midrule
\missing{phase4\_augmentation\_policy results} \\
\bottomrule
\end{tabular}
}\caption{Augmentation policies (pilot, validation).}\label{tab:policy}\end{table}

\subsection{Statistics}
The primary metric is the rare-class macro-F1, the mean of the S and F F1 scores. Uncertainty comes from a
cluster bootstrap over the 22 test records (10{,}000 replicates; seeds averaged within each replicate), because
beats of one patient are correlated. Paired differences resample the same records for both conditions;
two-sided percentile-bootstrap $p$-values are adjusted with the Holm procedure~\cite{holm1979} within the
confirmatory family (cVAE and WGAN-GP against none, for both architectures: 4 tests) and within the secondary
family (cVAE against WGAN-GP and the best generative condition against SMOTE, class-weighted loss and random
oversampling: 8 tests). Replicates in which S or F has no true beats are excluded and counted. We also report
per-class F1, one-vs-rest PR-AUC, AAMI sensitivity and positive predictivity for S and V, leave-one-record-out
estimates and results without record 202.

\section{Results}
Table~\ref{tab:hypotheses} summarises the verdict on each pre-registered hypothesis; the sections below give the
evidence.
\begin{table}[htbp]\centering\small\adjustbox{max width=\linewidth}{% Generated by python -m ecgaug.paper from the result files
\begin{tabular}{lp{5.2cm}lp{5.2cm}}
\toprule
Hypothesis & Statement & Verdict & Evidence \\
\midrule
H1 & confirmatory: a generative condition improves test rare-class macro-F1 over none (Holm, 4 tests) & \missing{pending} &  \\
H1b & exploratory: gains larger for F than for S & \missing{pending} &  \\
H1c & secondary: generative beats SMOTE, random oversampling and class-weighted loss & \missing{pending} &  \\
H2 & exploratory: intra-patient evaluation inflates the baseline and the gains & \missing{pending} &  \\
H3 & exploratory: a quality metric predicts the downstream gain & \missing{pending} &  \\
H4 & exploratory: augmented models are at least as faithful as the baseline & \missing{pending} &  \\
\bottomrule
\end{tabular}
}\caption{Pre-registered hypotheses and their verdicts, computed from the result files by the registered criteria.}\label{tab:hypotheses}\end{table}

\subsection{RQ1: does generative augmentation help?}
\missing{test results: written after /final-eval}

\begin{table}[htbp]\centering\small\adjustbox{max width=\linewidth}{\missing{test results: written by /final-eval}
}\caption{Test results per condition (95\% record-level CI).}\label{tab:test}\end{table}
\begin{table}[htbp]\centering\small\adjustbox{max width=\linewidth}{\missing{test results: written by /final-eval}
}\caption{Pre-registered comparisons (unadjusted 95\% CI, Holm-adjusted $p$).}\label{tab:comparisons}\end{table}

\subsection{RQ2: how much does intra-patient evaluation inflate results?}
\missing{leakage audit: optional, runs after the final evaluation}

\begin{table}[htbp]\centering\small\adjustbox{max width=\linewidth}{% Generated by python -m ecgaug.paper from: nothing yet
\begin{tabular}{l}
\toprule
Condition \\
\midrule
\missing{leakage audit (optional, after the final evaluation)} \\
\bottomrule
\end{tabular}
}\caption{Rare-class macro-F1 of the cnn when the same pipeline is evaluated on a random beat-level split of DS1 and DS2 (intra-patient) and on DS2 (inter-patient). Run after the final evaluation; never used for any decision.}\label{tab:audit}\end{table}

\subsection{RQ3: do quality metrics predict the benefit?}
\missing{RQ3: python -m ecgaug.eval.rq3 after the phase 3 runs}

\begin{table}[htbp]\centering\small\adjustbox{max width=\linewidth}{% Generated by python -m ecgaug.paper from: nothing yet
\begin{tabular}{l}
\toprule
Metric \\
\midrule
\missing{RQ3: python -m ecgaug.eval.rq3} \\
\bottomrule
\end{tabular}
}\caption{Quality metrics against the downstream gain across generator checkpoints (exploratory; TSTR analysed apart).}\label{tab:rq3}\end{table}

\subsection{RQ4: does augmentation change what the classifier relies on?}
\missing{RQ4: the phase 5 gate}

\begin{table}[htbp]\centering\small\adjustbox{max width=\linewidth}{% Generated by python -m ecgaug.paper from: nothing yet
\begin{tabular}{l}
\toprule
Method \\
\midrule
\missing{phase5\_explainability results} \\
\bottomrule
\end{tabular}
}\caption{Saliency methods: faithfulness gain over random deletion and stability across seeds.}\label{tab:explain}\end{table}

\section{Grounded clinical notes}
For a classified beat, the note generator receives only a JSON object: an anonymous record identifier, the
predicted class, the four class probabilities, a low-confidence flag (top probability below 0.6) and the two
segments (P, QRS, ST, T) with the most saliency mass. Four generators were compared: a deterministic template,
and three LLM variants (a constrained prompt, few-shot examples with a fixed output schema, and a fact-check and
regenerate loop that falls back to the template). The fact checker verifies that every class name, number and
segment in a note appears in the input, that the disclaimer is present, and that low-confidence notes ask for
clinician review; the hard gate is zero unsupported claims over 50 validation cases.
\missing{note generators: the phase 5 gate}

\begin{table}[htbp]\centering\small% Generated by python -m ecgaug.paper from: nothing yet
\begin{tabular}{l}
\toprule
Notes \\
\midrule
\missing{phase5\_clinical\_notes results} \\
\bottomrule
\end{tabular}
\caption{Note generators: unsupported claims and coverage.}\label{tab:notes}\end{table}

\section{Limitations and ethics}
\paragraph{Small test population.} DS2 holds 22 records, 21 from patients not seen in training and one (202) from
the subject of DS1 record 201; S and F beats come from few of them, so intervals are wide and results may reflect
individual patients. We resample records rather than beats for this reason and report per-record and
leave-one-record-out results.
\paragraph{One database, one lead.} MIT-BIH was recorded decades ago at one hospital with narrow demographics; we
use one lead and reference R-peak positions, so detector errors are not modelled, and we did not validate on an
external database.
\paragraph{Generators.} The generators learn minority classes from very few patients, so synthetic diversity is
bounded by them; RR features of synthetic beats are copied from real beats, so generative augmentation can only
add morphology, which matters most for classes defined by timing such as S.
\paragraph{Selection.} Selection relies on a small validation set of DS1 patients, with few F beats.
\paragraph{Explanations.} Saliency methods are approximate, and the segment windows are fixed R-relative
approximations of P, QRS, ST and T.
\paragraph{Notes and clinical use.} Notes were rated by team members, not clinicians, and restate a model that is
not a medical device; every note carries a disclaimer and is not intended for clinical decisions. The data are
public and de-identified; synthetic banks derived from them are released only if the PhysioNet licence allows it.

\section{Reproducibility}
Every run is defined by its configuration file, seed and git commit, and is logged with its environment,
training curve, validation predictions and checkpoint hash; \texttt{results/runs\_index.csv} maps every table cell
to run identifiers. The selection rules were committed before the experiments they govern
(\texttt{docs/PREREGISTRATION.md}), every gate decision is recorded with its evidence
(\texttt{docs/DECISIONS.md}), frozen artifacts carry SHA-256 hashes (\texttt{results/manifest.json}), and
\texttt{docs/TEST\_ACCESS\_LOG.md} shows the single evaluation on the test set. \texttt{scripts/reproduce.sh} runs the
pipeline in order, and a dress rehearsal on synthetic records exercises every stage. Code, configurations and the
lockfile are released with the paper.

\bibliographystyle{plain}
\bibliography{refs}
\end{document}
